\section{Introduction}
\label{sec:intro}

Deep learning has entered an era characterized by large-scale pre-trained foundation models that exhibit zero-shot capabilities across diverse domains \cite{radford2021learning, devlin2018bert, vaswani2017attention}. To adapt these general-purpose base models to specific downstream applications, practitioners commonly employ fine-tuning, generating multiple task-specific expert models \cite{wortsman2022robust}. However, deploying and maintaining a dedicated, full-sized neural network for every individual task quickly becomes computationally intensive. To address this bottleneck, researchers have turned to \emph{model merging}—the process of consolidating multiple task-specific networks into a single multi-task model without additional training or inference overhead \cite{fedavg, wortsman2022model, ainsworth2022git}.

One of the simplest and most widely used model merging frameworks is \emph{Task Arithmetic} (TA) \cite{ilharco2022editing}, which linearly combines task-specific weight updates (referred to as task vectors) and adds them back to the pre-trained base model. While Task Arithmetic is simple, parameter-free, and computationally efficient, it frequently suffers from \emph{destructive interference}—a phenomenon where overlapping parameter updates for different tasks cancel each other out, leading to representation degradation and a drop in multi-task performance \cite{yadav2023ties, yu2023dare}.

To mitigate this interference, recent literature has proposed increasingly complex and highly parameterized merging methodologies. For instance, FoldMerge (Neural Origami) \cite{foldmerge2025} introduces a non-linear RealNVP normalizing flow diffeomorphism featuring approximately $2.6\text{M}$ auxiliary parameters to warp weight coordinates into a shared space. Other frameworks like SyMerge \cite{symerge2025} and AdaMerging \cite{yang2023adamerging} formulate merging as an optimization problem, utilizing iterative test-time adaptation (TTA) on test streams or validation sets to search for layer-wise coefficients. 

While these complex pipelines can achieve strong merging performance, they introduce parameter and computational overhead. For example, optimizing normalizing flows or running test-time adaptation can require several minutes of optimization on high-performance hardware, which can be challenging in resource-constrained deployment environments. Second, as shown in concurrent studies \cite{saim2025}, these high-dimensional optimization objectives can suffer from the ``Overfitting-Optimizer Paradox,'' where optimization parameters overfit to specific evaluation streams, resulting in fragile representations that fail to generalize. Third, as we demonstrate, substantial multi-task consolidation can be achieved through closed-form linear algebra transformations without any iterative optimization.

Guided by Occam's razor, we ask: \emph{Can we reduce destructive parameter interference and merge model weights successfully without auxiliary parameters, optimization steps, or test-time adaptation?}

In this work, we propose \textbf{Spectral Model Merging via Singular Value Slicing (SVS)}, a non-parametric, training-free, and closed-form model merging framework. SVS is built on the hypothesis that the essential knowledge acquired during downstream fine-tuning resides in a low-dimensional spectral manifold of the weight updates, while the remaining high-dimensional components consist of sequential fine-tuning noise that causes destructive interference when merged.

To isolate this core knowledge, SVS performs a standard Singular Value Decomposition (SVD) on each task-specific delta matrix. It then applies a low-rank projection operator—which we refer to as \emph{Singular Value Slicing}—retaining only the top $k$ principal singular values and vectors before linearly combining the task vectors. By discarding the high-frequency spectral components, SVS acts as an analytical noise filter. 

Because slicing and combining parameters can distort the activation scale of the merged network, we also discuss \textbf{Barycentric Weight Normalization (BWN)}, an analytical scale-preservation operator designed to match the merged weights' Frobenius norm to the weighted barycenter of the experts. Crucially, we show mathematically that standard downstream feature normalization layers (such as LayerNorm, RMSNorm, or L2-normalization) mathematically neutralize global weight scaling factors. Consequently, BWN is primarily relevant in un-normalized neural network environments (such as standard multi-layer MLPs or convolutional layers without normalization). We provide a formal proof of this scale-invariance and empirically validate BWN's utility in un-normalized MLP environments.

We evaluate SVS on full-network model merging of all 86M parameters of the CLIP-ViT-B/32 vision encoder across four diverse image classification datasets: MNIST, FashionMNIST, CIFAR-10, and SVHN. Our empirical results demonstrate that:
\begin{itemize}
    \item \textbf{Full-Network Spectral Merging:} SVS successfully consolidates four specialized task experts into a single multi-task model by merging all 86 million parameters of the visual backbone. SVS robustly tracks Task Arithmetic within $1.55\%$ at $98.0\%$ spectral rank compression (rank $k=16$), while strictly matching or exceeding it at rank $k=128$ ($16.7\%$ of the available rank space).
    \item \textbf{Regularization in Intermediate Regimes:} In intermediate coefficient regimes (e.g., $\lambda \in [0.3, 0.5]$), SVS tracks or slightly exceeds standard Task Arithmetic, showing that low-rank slicing can act as an analytical regularizer.
    \item \textbf{Zero Test-Time Computational Overhead:} SVS requires no trainable parameters or test-time optimization, and runs analytically in less than a minute (approximately 55 seconds) for the entire 86M parameter backbone on a standard CPU node. Crucially, once the one-time Singular Value Decomposition is computed, subsequent merges (such as sweeping ranks or evaluating dynamic entropy-based allocations) are accelerated via SVD Caching to approximately 1.2 seconds (fully dominated by PyTorch model deep copies), presenting a highly efficient and scalable merging baseline.
\end{itemize}

By demonstrating that model merging can be achieved using closed-form linear algebra, SVS serves as an efficient and simple baseline, encouraging further research into non-parametric, spectral-filtering merging methods.
